\hypertarget{mux1ee5c-tiuxeau-tuux1ea7n-2}{%
\subsubsection{Mục tiêu tuần 2:}}

\begin{itemize}
\tightlist
\item
  Nghiên cứu sâu về Scrapy framework --- công cụ chính để crawl tin tức.
\item
  Tìm hiểu sự khác biệt giữa CrawlSpider (v1) và SitemapSpider (v2).
\item
  Nắm vững Docker, Dockerfile, docker-compose để chuẩn bị containerize.
\item
  Phân tích cấu trúc sitemap XML của 3 báo điện tử Việt Nam.
\end{itemize}

\hypertarget{cuxe1c-cuxf4ng-viux1ec7c-cux1ea7n-triux1ec3n-khai-trong-tuux1ea7n-nuxe0y}{%
\subsubsection{Các công việc cần triển khai trong tuần
này:}}

\begingroup
\small
\setlength{\tabcolsep}{3pt}
\renewcommand{\arraystretch}{1.15}
\begin{longtable}[]{@{}
  >{\raggedright\arraybackslash}p{(\columnwidth - 8\tabcolsep) * \real{0.0121}}
  >{\raggedright\arraybackslash}p{(\columnwidth - 8\tabcolsep) * \real{0.7126}}
  >{\raggedright\arraybackslash}p{(\columnwidth - 8\tabcolsep) * \real{0.0486}}
  >{\raggedright\arraybackslash}p{(\columnwidth - 8\tabcolsep) * \real{0.0607}}
  >{\raggedright\arraybackslash}p{(\columnwidth - 8\tabcolsep) * \real{0.1660}}@{}}
\toprule\noalign{}
\begin{minipage}[b]{\linewidth}\raggedright
Thứ
\end{minipage} & \begin{minipage}[b]{\linewidth}\raggedright
Công việc
\end{minipage} & \begin{minipage}[b]{\linewidth}\raggedright
Ngày bắt đầu
\end{minipage} & \begin{minipage}[b]{\linewidth}\raggedright
Ngày hoàn thành
\end{minipage} & \begin{minipage}[b]{\linewidth}\raggedright
Nguồn tài liệu
\end{minipage} \\
\midrule\noalign{}
\endhead
\bottomrule\noalign{}
\endlastfoot
2 & - Tìm hiểu Scrapy framework: \qquad  + Cấu trúc project Scrapy
\qquad  + Spider, Pipeline, Settings \qquad  + Middleware và Item
Pipeline & 09/06/2026 & 09/06/2026 & \url{https://docs.scrapy.org/} \\
3 & - So sánh CrawlSpider vs SitemapSpider: \qquad  + CrawlSpider: dựa
vào DEPTH\_LIMIT, bỏ sót bài cũ \qquad  + SitemapSpider: đọc toàn bộ URL
từ sitemap, không phụ thuộc navigation & 10/06/2026 & 10/06/2026 &
\url{https://docs.scrapy.org/} \\
4 & - Phân tích cấu trúc sitemap XML của 3 báo: \qquad  + VnExpress:
\texttt{sitemap\_news.xml} \qquad  + Thanh Niên: \texttt{sitemap.xml}
\qquad  + VietnamNet: \texttt{sitemap\_news.xml} & 11/06/2026 &
11/06/2026 & \\
5 & - Tìm hiểu Docker cơ bản: \qquad  + Dockerfile: FROM, WORKDIR, COPY,
RUN, CMD \qquad  + docker-compose.yml: services, volumes, environment
\qquad  + Multi-stage build & 12/06/2026 & 12/06/2026 &
\url{https://docs.docker.com/} \\
6 & - Tìm hiểu thư viện \texttt{newspaper3k} để extract article content
- Nghiên cứu Kafka Producer/Consumer pattern - Tìm hiểu
\texttt{confluent\_kafka} Python client & 13/06/2026 & 13/06/2026 &
\url{https://newspaper.readthedocs.io/} \\
\end{longtable}
\endgroup

\hypertarget{kux1ebft-quux1ea3-ux111ux1ea1t-ux111ux1b0ux1ee3c-tuux1ea7n-2}{%
\subsubsection{Kết quả đạt được tuần
2:}}

\begin{itemize}
\tightlist
\item
  Nắm vững cấu trúc project Scrapy:

  \begin{itemize}
  \tightlist
  \item
    \texttt{spiders/}: chứa các spider class
  \item
    \texttt{pipelines.py}: xử lý item sau khi spider yield
  \item
    \texttt{settings.py}: cấu hình CONCURRENT\_REQUESTS,
    DOWNLOAD\_DELAY, DEPTH\_LIMIT,\ldots{}
  \end{itemize}
\item
  Hiểu rõ sự khác biệt giữa 2 loại Spider:

  \begin{itemize}
  \tightlist
  \item
    \textbf{CrawlSpider (v1)}: Dùng \texttt{DEPTH\_LIMIT=5}, chỉ crawl
    được trang trong phạm vi liên kết \(\rightarrow\) bỏ sót bài cũ
    không được link từ trang chủ
  \item
    \textbf{SitemapSpider (v2)}: Đọc trực tiếp từ file sitemap XML
    \(\rightarrow\) crawl toàn bộ bài viết bao gồm cả bài cũ, không bị
    timeout 15 phút như Lambda
  \end{itemize}
\item
  Đã phân tích sitemap XML của 3 báo:

  \begin{itemize}
  \tightlist
  \item
    VnExpress: \texttt{https://vnexpress.net/sitemap\_news.xml}
  \item
    Thanh Niên: \texttt{https://thanhnien.vn/sitemap.xml}
  \item
    VietnamNet: \texttt{https://vietnamnet.vn/sitemap\_news.xml}
  \end{itemize}
\item
  Nắm được kiến thức Docker cần thiết:

  \begin{itemize}
  \tightlist
  \item
    Cách viết Dockerfile tối ưu (base image \texttt{python:3.10-slim},
    multi-layer caching)
  \item
    Docker Compose để chạy PostgreSQL + Kafka local
  \item
    Hiểu biến môi trường \texttt{.env} và cách inject vào container
  \end{itemize}
\item
  Tìm hiểu \texttt{newspaper3k}:

  \begin{itemize}
  \tightlist
  \item
    Thư viện Python hỗ trợ extract title, content, author, publish\_date
    từ HTML
  \item
    Dùng \texttt{article.set\_html()} + \texttt{article.parse()} thay vì
    download trực tiếp
  \item
    Phối hợp với CSS Selectors của Scrapy để tăng độ chính xác
  \end{itemize}
\end{itemize}
